% unit: u13
% sources: sources/065-graphs-networks-incidence-matrices.txt, sources/167-graphs-networks-incidence-matrices.txt, sources/168-mit18-06scf11-ses1-12sum-pdf-download-file.txt, sources/169-graphs-and-networks.txt, sources/170-mit18-06scf11-ses1-12prob-pdf-download-file.txt, sources/171-mit18-06scf11-ses1-12sol-pdf-download-file.txt
\section{Graphs, networks, incidence matrices}\label{sec:u13}

The matrices in the previous sections were invented on the spot. Matrices in applications come from
somewhere and have structure, and whoever works with them uses that structure. The most important
example in applied mathematics is the matrix of a graph. Here every one of the four subspaces has a
physical meaning (potentials, potential differences, currents, and Kirchhoff's current law), and
the dimension count of Section~\ref{sec:u11} proves Euler's formula for graphs.

\subsection{The incidence matrix of a directed graph}

A graph has nodes and edges (Definition~\ref{def:u12-graph}). To measure currents we give each edge a
direction, an arrow that says which way of flowing counts as positive; the choice of arrows is
arbitrary.

\begin{definition}[Incidence matrix]\label{def:u13-incidence}
Let a directed graph have $n$ nodes and $m$ edges. Its \emph{incidence matrix} $A$ is $m\times n$,
with one row per edge and one column per node. If edge $k$ goes from node $i$ to node $j$, row $k$
has $-1$ in column $i$, $+1$ in column $j$, and zeros elsewhere.
\end{definition}

\begin{example}[Four nodes, five edges]\label{ex:u13-graph}
The graph of the lecture (Figure~\ref{fig:u13-graph}, left) has $n=4$ nodes and $m=5$ edges:
edge~1 from node~1 to~2, edge~2 from~2 to~3, edge~3 from~1 to~3, edge~4 from~1 to~4, and edge~5
from~3 to~4. Its incidence matrix is
\[
A=\begin{bmatrix}-1&1&0&0\\0&-1&1&0\\-1&0&1&0\\-1&0&0&1\\0&0&-1&1\end{bmatrix}
\begin{matrix}\text{edge }1\\\text{edge }2\\\text{edge }3\\\text{edge }4\\\text{edge }5\end{matrix}.
\]
Each row has exactly two nonzeros, so a large graph gives a large but very sparse matrix ($2m$
nonzero entries out of $mn$). Edges $1,2,3$ form a loop, and the corresponding rows are dependent:
row~1 $+$ row~2 $=(-1,0,1,0)=$ row~3. Loops will turn out to be exactly the source of dependent
rows.
\end{example}

\begin{figure}[ht]\centering
\begin{tikzpicture}[>=Stealth,font=\small,
    nd/.style={circle,draw,fill=white,inner sep=1.5pt,minimum size=16pt},
    el/.style={fill=white,inner sep=1pt,font=\footnotesize}]
  \begin{scope}
    \node[nd] (a) at (0,0) {1};
    \node[nd] (b) at (2.2,1.6) {2};
    \node[nd] (c) at (4.4,0) {3};
    \node[nd] (d) at (2.2,-1.6) {4};
    \draw[->,line width=1.6pt,blue!70!black] (a) -- node[el] {1} (b);
    \draw[->,line width=1.6pt,blue!70!black] (b) -- node[el] {2} (c);
    \draw[->,thick] (a) -- node[el] {3} (c);
    \draw[->,line width=1.6pt,blue!70!black] (a) -- node[el] {4} (d);
    \draw[->,thick] (c) -- node[el] {5} (d);
    \draw[->,gray] (2.2,0.75) ++(250:0.32) arc[start angle=250,end angle=-30,radius=0.32];
    \draw[->,gray] (2.2,-0.75) ++(30:0.32) arc[start angle=30,end angle=-250,radius=0.32];
    \node at (2.2,-2.5) {edges $1,2,4$: a spanning tree};
  \end{scope}
  \begin{scope}[xshift=7cm]
    \node[nd,label={[font=\footnotesize]left:$x_1=\tfrac58$}] (a) at (0,0) {1};
    \node[nd,label={[font=\footnotesize]above:$x_2=\tfrac12$}] (b) at (2.2,1.6) {2};
    \node[nd,label={[font=\footnotesize]right:$x_3=\tfrac38$}] (c) at (4.4,0) {3};
    \node[nd,label={[font=\footnotesize]below:$x_4=0$}] (d) at (2.2,-1.6) {4};
    \draw[->,thick] (a) -- node[el] {$\tfrac18$} (b);
    \draw[->,thick] (b) -- node[el] {$\tfrac18$} (c);
    \draw[->,thick] (a) -- node[el] {$\tfrac14$} (c);
    \draw[->,thick] (a) -- node[el] {$\tfrac58$} (d);
    \draw[->,thick] (c) -- node[el] {$\tfrac38$} (d);
    \draw[->,very thick,orange!80!black] (-1.3,-0.9) -- (a) node[pos=0,below,font=\footnotesize] {$f_1=1$ in};
    \draw[->,very thick,orange!80!black] (d) -- (3.6,-2.3) node[right,font=\footnotesize] {$1$ out};
    \node at (2.2,-2.5) {};
  \end{scope}
\end{tikzpicture}
\caption{Left: the directed graph of Example~\ref{ex:u13-graph}; the two circular arrows mark
the independent loops (edges $1,2,3$ and edges $3,5,4$), and the thick edges $1,2,4$ form a tree.
Right: with unit conductances, one unit of current entering node~1 and leaving node~4
(Example~\ref{ex:u13-network}): node potentials and edge currents, all flowing along the arrows.}
\label{fig:u13-graph}
\end{figure}

\subsection{Potentials and the nullspace}

Multiply $A$ by a vector $\vb{x}=(x_1,x_2,x_3,x_4)$ of \emph{potentials} at the nodes:
\[
A\vb{x}=\begin{bmatrix}x_2-x_1\\x_3-x_2\\x_3-x_1\\x_4-x_1\\x_4-x_3\end{bmatrix}.
\]
Each component is a \emph{potential difference} across an edge (head minus tail). Potential
differences are what make current flow.

\begin{proposition}[Nullspace of an incidence matrix]\label{prop:u13-null}
If the graph is connected (every two nodes are joined by a path), then $\Nul{A}$ is the line of
constant vectors $c(1,1,\dots,1)$, and $\rank A=n-1$.
\end{proposition}
\begin{proof}
A constant vector gives zero differences, so it lies in $\Nul{A}$. Conversely, $A\vb{x}=\vb{0}$
says $x_i=x_j$ whenever an edge joins $i$ and $j$; following a path from node~1 to any node $k$
gives $x_k=x_1$. So $\dim\Nul{A}=1$ and $r=n-1$.
\end{proof}

Physically, potentials are determined only up to an additive constant, like the constant of
integration in calculus or the choice of zero on a temperature scale. One fixes it by
\emph{grounding} a node, say $x_4=0$. That removes column~4, and the remaining three columns are
independent. In the example $r=3$, and any three of the four columns are independent.

Which vectors $\vb{e}$ of edge values are potential differences, i.e.\ lie in $\Col{A}$? Going around
a loop, the differences telescope: on the loop of edges $1,2,3$,
$(x_2-x_1)+(x_3-x_2)-(x_3-x_1)=0$. So every $\vb{e}=A\vb{x}$ satisfies $e_1+e_2-e_3=0$, and
likewise $e_3+e_5-e_4=0$ on the other loop. This is \emph{Kirchhoff's voltage law}, and by
Theorem~\ref{thm:u09-solvable} these loop conditions are also sufficient once we know that the loop
vectors span the left nullspace, which we prove next.

\subsection{Currents and Kirchhoff's current law}

Now let $\vb{y}=(y_1,\dots,y_5)$ be \emph{currents} on the edges, positive when flowing along the
arrow. Then
\[
A\T\vb{y}=\begin{bmatrix}-1&0&-1&-1&0\\1&-1&0&0&0\\0&1&1&0&-1\\0&0&0&1&1\end{bmatrix}
\begin{bmatrix}y_1\\y_2\\y_3\\y_4\\y_5\end{bmatrix}
=\begin{bmatrix}-y_1-y_3-y_4\\y_1-y_2\\y_2+y_3-y_5\\y_4+y_5\end{bmatrix}.
\]
Row $i$ of $A\T\vb{y}$ is (current into node $i$) minus (current out of node $i$). The equation
\[
A\T\vb{y}=\vb{0}\qquad\text{(Kirchhoff's current law)}
\]
says that at every node, in equals out: charge does not accumulate. It is a balance equation, a
conservation law, perhaps the most fundamental equation of applied mathematics. Its solutions form
the left nullspace, of dimension $m-r=5-3=2$.

\begin{definition}[Loop vector]\label{def:u13-loop}
For a loop in the graph traversed in a chosen direction, its \emph{loop vector} has entry $+1$ for
each edge traversed along its arrow, $-1$ for each edge traversed against its arrow, and $0$ for
edges not on the loop.
\end{definition}

\begin{theorem}[Loops span the left nullspace]\label{thm:u13-loops}
Every loop vector satisfies $A\T\vb{y}=\vb{0}$, and every solution of $A\T\vb{y}=\vb{0}$ is a
combination of loop vectors.
\end{theorem}
\begin{proof}
A loop vector sends one unit of current around the loop. At each node on the loop, one unit
arrives along one loop edge and one unit leaves along the next, so in minus out is $0$; nodes off
the loop see no current.

Conversely, let $\vb{y}\ne\vb{0}$ satisfy the current law. Pick an edge with $y_k\ne0$ and walk along
it in the direction the current flows. At the node reached, current arrives, so by the balance
some other edge carries current away; walk along it in its direction of flow. Continue. With
finitely many nodes, the walk eventually revisits a node, closing a loop on which every edge carries
nonzero current in the walking direction. Let $t>0$ be the smallest of these currents in absolute
value and subtract $t$ times this loop's vector from $\vb{y}$. The result still satisfies the
current law and has at least one more zero entry. Repeating, after finitely many steps nothing is
left, so $\vb{y}$ was a sum of multiples of loop vectors.
\end{proof}

\begin{example}[Two independent loops]\label{ex:u13-loops}
Send one unit around the upper triangle $1\to2\to3\to1$: edges $1$ and $2$ are traversed along
their arrows and edge~$3$ against it, giving $\vb{y}=(1,1,-1,0,0)$. Check the four node equations:
$-1+1-0=0$, $1-1=0$, $1-1-0=0$, $0+0=0$. The lower triangle $1\to3\to4\to1$ gives
$(0,0,1,-1,1)$. These two are independent (look at entries $1$ and $5$), and $\dim\Nul{A\T}=2$,
so they form a basis. The big outer loop $1\to2\to3\to4\to1$ gives $(1,1,0,-1,1)$, which is
their sum: on edge~3 the two loop currents cancel. It is a loop, but not an independent one.
\end{example}

\subsection{Trees, the row space and Euler's formula}

The row space $\Col{A\T}$ has dimension $r=3$. Which three rows (edges) are independent? Row
reduction of $A\T$ finds pivots in columns $1,2,4$, the columns of $A\T$ that belong to edges $1,2$
and~$4$. Column~3 of $A\T$ is column~1 plus column~2, which is the loop $1,2,3$ again. The pivot edges
$1,2,4$ connect all four nodes and contain no loop (thick edges in Figure~\ref{fig:u13-graph}).
A graph with no loops is a \emph{tree}. Independent rows correspond to edges without loops;
dependencies among rows are exactly loops.

\begin{theorem}[Euler's formula]\label{thm:u13-euler}
For a connected graph with $n$ nodes, $m$ edges and $\ell=\dim\Nul{A\T}$ independent loops,
\[
n-m+\ell=1.
\]
\end{theorem}
\begin{proof}
By Theorem~\ref{thm:u11-dims}, $\ell=\dim\Nul{A\T}=m-r$, and $r=n-1$ by
Proposition~\ref{prop:u13-null}. So $\ell=m-n+1$.
\end{proof}

Nodes are zero-dimensional, edges one-dimensional, and loops enclose two-dimensional regions; the
alternating count is always $1$. For the example, $4-5+2=1$. For a graph with $5$ nodes and $7$
edges drawn in the plane with $3$ independent loops, $5-7+3=1$. Linear algebra has proved a theorem
of topology.

\begin{trap}
Euler's formula counts \emph{independent} loops. The graph of Figure~\ref{fig:u13-graph} has three
visible loops (two triangles and the outer quadrilateral), and $4-5+3\neq1$; the outer loop is the
sum of the other two. Also, Proposition~\ref{prop:u13-null} needs a connected graph: if the graph
falls into separate pieces, $\Nul{A}$ is larger (Exercise~\ref{exr:u13-disconnected}).
\end{trap}

\subsection{\texorpdfstring{Ohm's law and \(A\T CA\vb{x}=\vb{f}\)}{Ohm's law and ATCAx = f}}

The three steps of an electrical network chain together:
\begin{center}
\begin{tikzpicture}[>=Stealth,font=\small,node distance=2.6cm]
  \node[align=center] (x) {$\vb{x}$\\potentials\\at nodes};
  \node[align=center,right=of x] (e) {$\vb{e}$\\potential\\differences};
  \node[align=center,right=of e] (y) {$\vb{y}$\\currents\\on edges};
  \node[align=center,right=of y] (f) {$\vb{f}$\\current sources\\at nodes};
  \draw[->,thick] (x) -- node[above] {$\vb{e}=A\vb{x}$} (e);
  \draw[->,thick] (e) -- node[above] {$\vb{y}=-C\vb{e}$} node[below,font=\footnotesize] {Ohm} (y);
  \draw[->,thick] (y) -- node[above] {$A\T\vb{y}=-\vb{f}$} node[below,font=\footnotesize] {Kirchhoff} (f);
\end{tikzpicture}
\end{center}
Ohm's law says that the current on an edge is its conductance (one over its resistance) times the
potential drop; $C$ is the diagonal matrix of conductances. Current flows from high to low
potential, hence the minus sign in $\vb{y}=-C\vb{e}$. If a current $f_i$ is fed into node $i$ from
outside, the balance at that node becomes (in) $-$ (out) $+f_i=0$, i.e.\ $A\T\vb{y}=-\vb{f}$.
Combining the three equations, the signs cancel:
\[
A\T CA\,\vb{x}=\vb{f}.
\]
(The lecture writes the chain as $\vb{e}=A\vb{x}$, $\vb{y}=C\vb{e}$, $A\T\vb{y}=\vb{f}$, absorbing
the signs into the conventions; the resulting equation is the same.) This is the basic equation of
equilibrium for networks, and the matrix $A\T CA$ is symmetric:
$(A\T CA)\T=A\T C\T A=A\T CA$.

\begin{proposition}[The matrix \texorpdfstring{\(A\T A\)}{ATA} of a graph]\label{prop:u13-laplacian}
For an incidence matrix $A$, the diagonal entry $(A\T A)_{ii}$ is the \emph{degree} of node $i$
(the number of edges touching it), the off-diagonal entry $(A\T A)_{ij}$ is $-1$ if an edge joins
$i$ and $j$ and $0$ otherwise, and $\tr(A\T A)=2m$.
\end{proposition}
\begin{proof}
$(A\T A)_{ij}$ is the dot product of columns $i$ and $j$ of $A$. Column $i$ has an entry $\pm1$ in
each row whose edge touches node $i$, so its squared length is the degree. Two different columns
share a nonzero row only for an edge joining $i$ and $j$, where the entries are $-1$ and $+1$. The
trace is the sum of the degrees, and each edge contributes to two degrees.
\end{proof}

\begin{example}[Solving a network]\label{ex:u13-network}
Take all conductances equal to $1$ ($C=I$) in the graph of Example~\ref{ex:u13-graph}, feed one unit
of current into node~1 and take it out at node~4: $\vb{f}=(1,0,0,-1)$. The degrees are $3,2,3,2$, so
\[
A\T A=\begin{bmatrix}3&-1&-1&-1\\-1&2&-1&0\\-1&-1&3&-1\\-1&0&-1&2\end{bmatrix}.
\]
It is singular (its rows add to zero, reflecting $\Nul{A}$), so ground node~4: set $x_4=0$ and
drop the fourth equation. The remaining equations $3x_1-x_2-x_3=1$, $-x_1+2x_2-x_3=0$,
$-x_1-x_2+3x_3=0$ give, by elimination, $\vb{x}=(\tfrac58,\tfrac12,\tfrac38,0)$. Check the first:
$\tfrac{15}{8}-\tfrac48-\tfrac38=1$. The potential differences and currents are
\[
\vb{e}=A\vb{x}=\bigl(-\tfrac18,-\tfrac18,-\tfrac14,-\tfrac58,-\tfrac38\bigr),\qquad
\vb{y}=-\vb{e}=\bigl(\tfrac18,\tfrac18,\tfrac14,\tfrac58,\tfrac38\bigr).
\]
All currents are positive, so all flow along the arrows (Figure~\ref{fig:u13-graph}, right). Node~1
sends out $\tfrac18+\tfrac14+\tfrac58=1$, node~2 passes $\tfrac18$ through, node~3 receives
$\tfrac18+\tfrac14$ and sends $\tfrac38$, and node~4 collects $\tfrac58+\tfrac38=1$. The direct edge
$1\to4$ carries the most current, because every other route to node~4 crosses two edges.
\end{example}

\subsection{Exercises}

\begin{exercise}\label{exr:u13-square}
A square graph has nodes $1,2,3,4$ and edges $1\to2$, $2\to3$, $4\to3$, $1\to4$ (in that order).
Write its $4\times4$ incidence matrix $A$. Which vectors are in $\Nul{A}$? Why is $(1,0,0,0)$ not in
the row space of $A$?
\end{exercise}
\begin{solution}
\[
A=\begin{bmatrix}-1&1&0&0\\0&-1&1&0\\0&0&1&-1\\-1&0&0&1\end{bmatrix},\qquad
A\vb{x}=\begin{bmatrix}x_2-x_1\\x_3-x_2\\x_3-x_4\\x_4-x_1\end{bmatrix}.
\]
$A\vb{x}=\vb{0}$ forces $x_1=x_2=x_3=x_4$, so $\Nul{A}$ is all $(a,a,a,a)$. Every row of $A$ has
entries adding to $0$ (one $-1$ and one $+1$), so every combination of rows has entries adding to
$0$; the entries of $(1,0,0,0)$ add to $1$, so it is not in the row space.
\end{solution}

\begin{exercise}\label{exr:u13-ATCA}
For the square graph of Exercise~\ref{exr:u13-square} let the conductances be
$C=\operatorname{diag}(1,2,2,1)$. Find $A\T CA$. For the current source $\vb{f}=(1,0,-1,0)$ (one
unit into node~1, out of node~3), solve $A\T CA\vb{x}=\vb{f}$ with node~3 grounded, and find the
currents $\vb{y}=-CA\vb{x}$.
\end{exercise}
\begin{solution}
Multiplying out (or using Proposition~\ref{prop:u13-laplacian} with each edge weighted by its
conductance),
\[
A\T CA=\begin{bmatrix}2&-1&0&-1\\-1&3&-2&0\\0&-2&4&-2\\-1&0&-2&3\end{bmatrix}.
\]
With $x_3=0$, the equations from nodes $1,2,4$ are $2x_1-x_2-x_4=1$, $-x_1+3x_2=0$,
$-x_1+3x_4=0$. So $x_2=x_4=x_1/3$ and $2x_1-\tfrac23x_1=1$, giving $x_1=\tfrac34$ and
$\vb{x}=(\tfrac34,\tfrac14,0,\tfrac14)$. Then $A\vb{x}=(-\tfrac12,-\tfrac14,-\tfrac14,-\tfrac12)$ and
$\vb{y}=-CA\vb{x}=(\tfrac12,\tfrac12,\tfrac12,\tfrac12)$: half the current takes the route
$1\to2\to3$ and half takes $1\to4\to3$, as symmetry demands (each route has conductances $1$ and $2$
in series). Check at node~1: $\tfrac12+\tfrac12=1$ leaves.
\end{solution}

\begin{exercise}\label{exr:u13-five}
A graph has $5$ nodes and $6$ edges: $1\to2$, $2\to3$, $1\to3$, $2\to4$, $4\to5$, $5\to3$. Without
row reduction, find $\Nul{A}$, a basis of $\Nul{A\T}$, and $\tr(A\T A)$.
\end{exercise}
\begin{solution}
The graph is connected, so $\Nul{A}$ is the line through $(1,1,1,1,1)$ and $r=4$; then
$\dim\Nul{A\T}=6-4=2$. Loop vectors: the triangle $1\to2\to3\to1$ gives $(1,1,-1,0,0,0)$, and the
square $2\to4\to5\to3\to2$ gives $(0,-1,0,1,1,1)$ (edge~2 is traversed against its arrow). They are
independent, so they form a basis; the outer loop $1\to2\to4\to5\to3\to1$ gives their sum
$(1,0,-1,1,1,1)$. The degrees are $2,3,3,2,2$, so $\tr(A\T A)=12=2m$.
\end{solution}

\begin{exercise}\label{exr:u13-disconnected}
A graph consists of two separate triangles: nodes $1,2,3$ with edges $1\to2$, $2\to3$, $1\to3$, and
nodes $4,5,6$ with edges $4\to5$, $5\to6$, $4\to6$. Find $\dim\Nul{A}$, the rank, and
$\dim\Nul{A\T}$. What becomes of Euler's formula?
\end{exercise}
\begin{solution}
$A\vb{x}=\vb{0}$ forces equal potentials within each triangle but not across, so $\Nul{A}$ has basis
$(1,1,1,0,0,0)$ and $(0,0,0,1,1,1)$: dimension $2$, rank $6-2=4$. Then
$\dim\Nul{A\T}=m-r=6-4=2$, one loop per triangle. Now $n-m+\ell=6-6+2=2$: in general the alternating
count equals the number of connected pieces, and Euler's formula is the case of one piece.
\end{solution}
